\answer{ex:04:1}{(a)~$\approx1.1\cdot10^{11}$ bits/J. (b)~$\log_2\sqrt{10}\approx1.7$ bits frente a $81$: $\approx50$ veces menos.}
\answer{ex:04:2}{$r^*=(P_0/E_{\text{esp}})\ln\bigl(1/(r^*\Delta t)\bigr)$, $P_0/E_{\text{esp}}=1.16\,\text{s}^{-1}$: $r^*\approx6$ espigas por segundo, $7.4\cdot10^{10}$ bits/J.}
\answer{ex:04:3}{$\mathcal I=\frac{N}{1+(N-1)c}=9.9$ (límite $10$) frente a $\mathcal I=\frac{N}{1-c}=1111$: $\approx110$ veces; la correlación solo es dañina cuando se parece a la señal.}
\answer{ex:04:4}{$\theta\approx68.7$ y $19.9$ (en unidades de $w$), $m=19$ y $m=10$: el receptor rápido necesita casi la mitad de entradas coincidentes, aunque su entrada media es cinco veces menor.}
\answer{ex:04:5}{$U\tau_{\text{rec}}\ge0.725\,\text{s}$ y $\tau_{\text{rec}}(1-2U)\le0.05\,\text{s}$, lo que exige $U\ge0.483$; el menor $\tau_{\text{rec}}=0.725\,\text{s}$ con $U=1$ ($1.45\,\text{s}$ con $U=0.5$).}
\answer{ex:04:6}{(a)~$\theta\,e/(e-1)\approx1.58\,\theta$ para cualesquiera $\tau$ y $\theta$. (b)~$14$ espigas.}
\answer{ex:04:7}{(a)~$1.62$ bits por palabra: $0.77$ del número de espigas y $0.84$ de su disposición. (b)~$1.13$ y $1.59$ bits: con $30$ palabras la estimación pierde casi un tercio.}
